% Αναφορά διορθώσεων και ελέγχων. Μεταγλώττιση: latexmk -xelatex report.tex
% (τα results.tex και timings.tex παράγονται από το docs/build_docs.py)
\documentclass[a4paper,11pt]{article}
\usepackage{fontspec}
\usepackage[english,greek]{babel}
\setmainfont{LibertinusSerif}[Extension=.otf, UprightFont=*-Regular,
	BoldFont=*-Bold, ItalicFont=*-Italic, BoldItalicFont=*-BoldItalic]
\setsansfont{LibertinusSans}[Extension=.otf, UprightFont=*-Regular,
	BoldFont=*-Bold, ItalicFont=*-Italic]
\setmonofont{NewCMMono10}[Extension=.otf, UprightFont=*-Regular,
	BoldFont=*-Bold, Scale=MatchLowercase]
\usepackage[margin=2.2cm]{geometry}
\usepackage{graphicx,booktabs,longtable,array,xcolor,enumitem}
\usepackage[hidelinks]{hyperref}
\definecolor{ok}{HTML}{1F7A3A}
\definecolor{bad}{HTML}{B3261E}
\definecolor{frame}{HTML}{B9B2A6}
\newcommand{\PASS}{\textcolor{ok}{\textbf{PASS}}}
\newcommand{\FAIL}{\textcolor{bad}{\textbf{FAIL}}}
\newcommand{\shot}[2][\linewidth]{%
	{\setlength{\fboxsep}{0pt}\color{frame}\fbox{\includegraphics[width=#1]{#2}}}}
\setlist{itemsep=2pt,topsep=4pt}
\newcommand{\enc}{\texttt{\textbackslash en\{\}}}
% caption + screenshot kept on the same page
\newcommand{\evidence}[2]{\par\medskip\noindent
	\begin{minipage}{\linewidth}#1\par\smallskip\shot{#2}\end{minipage}\par}
\graphicspath{{../img/}}
\title{\textbf{Πρότυπο Διπλωματικής Εργασίας ΣΗΜΜΥ ΕΜΠ}\\[4pt]
	\large Αναφορά διορθώσεων και απόδειξη λειτουργίας}
\newcommand{\testcount}{19}
\newcommand{\testpassed}{19}
\newcommand{\pagesample}{80}
\newcommand{\testdate}{2026-10-11 01:20}
\newcommand{\texversion}{pdfTeX 3.141592653-2.6-1.40.25 (TeX Live 2023/Debian)}
\newcommand{\origpdferrors}{76}
\newcommand{\origxeerrors}{116}
\date{Έλεγχοι: 2026-10-11 01:20 · pdfTeX 3.141592653-2.6-1.40.25}

\begin{document}
\maketitle

\section{Σύνοψη}
Το template \texttt{estamos/NTUA-ECE-Thesis-Template} μεταγλωττιζόταν με
pdfLaTeX μόνο με την επιλογή \texttt{plain}, που απενεργοποιεί τα γραφικά.
Χωρίς αυτήν, με το υδατογράφημα του \texttt{main.tex} και τα σημερινά TeX Live
(και στο Overleaf), έβγαζε δεκάδες σφάλματα TikZ.
Με XeLaTeX δεν μεταγλωττιζόταν καθόλου. Επιπλέον, το αγγλικό κείμενο έπρεπε
να μπαίνει σε \verb|\en{...}|, αλλιώς τυπωνόταν με ελληνικά γράμματα.

Η νέα έκδοση:
\begin{itemize}
	\item μεταγλωττίζεται χωρίς σφάλματα με \textbf{pdfLaTeX, XeLaTeX και LuaLaTeX},
		με όλους τους συνδυασμούς επιλογών·
	\item με XeLaTeX/LuaLaTeX δέχεται \textbf{ελληνικά και αγγλικά μαζί χωρίς}
		\verb|\en{}| και οποιαδήποτε γραμματοσειρά OpenType·
	\item διορθώνει προϋπάρχοντα σφάλματα του pdfLaTeX (κώδικας, URL,
		ερωτηματικό, βιβλιογραφία)·
	\item συνοδεύεται από \textbf{\testcount{} αυτόματους ελέγχους}·
		πέρασαν \testpassed{} από \testcount{} (ενότητα~\ref{sec:tests}).
\end{itemize}
Τα παλιά έγγραφα με \verb|\en{...}| συνεχίζουν να δουλεύουν χωρίς αλλαγές.

\section{Προβλήματα, αιτίες και διορθώσεις}
\begin{longtable}{>{\raggedright}p{0.30\linewidth}>{\raggedright}p{0.30\linewidth}>{\raggedright\arraybackslash}p{0.32\linewidth}}
\toprule
\textbf{Πρόβλημα} & \textbf{Αιτία} & \textbf{Διόρθωση}\\
\midrule\endhead
Χωρίς \texttt{plain}, με \texttt{watermark}: σφάλματα «\foreignlanguage{english}{Giving up on this path. Did you forget a semicolon?}» σε κάθε σελίδα
& Το υδατογράφημα (TikZ) άλλαζε γλώσσα σε ελληνικά μέσα στο σχέδιο· το νεότερο babel-greek κάνει το \texttt{;} ενεργό
& Το υδατογράφημα γράφεται σε Unicode χωρίς αλλαγή γλώσσας· φορτώνεται η βιβλιοθήκη \texttt{babel} του TikZ\\
Χρονικό όριο στο Overleaf
& Δωρεάν Overleaf: όριο 10\,s από 9/2025· το πλήρες template θέλει 3–4 περάσματα
& Το λογότυπο δίνεται έτοιμο σε PDF (όχι μετατροπή EPS σε κάθε compile)· οδηγίες για \texttt{plain,noindex}\\
XeLaTeX/LuaLaTeX: \origxeerrors{} σφάλματα
& Η κλάση φόρτωνε \texttt{inputenc}/LGR και γραμματοσειρές Type~1
& Ανίχνευση μηχανής· \texttt{fontspec} με Libertinus ή New Computer Modern\\
Αγγλικό κείμενο χωρίς \verb|\en{}| τυπώνεται ελληνικά
& Κωδικοποίηση LGR του pdfLaTeX
& Με XeLaTeX/LuaLaTeX δεν χρειάζεται πλέον\\
Κώδικας \texttt{lstlisting}: \texttt{def}~→~«δεφ», χάνονται τα ελληνικά σχόλια
& Το listings τύπωνε με LGR· τα UTF-8 ελληνικά σπάνε σε bytes
& Λατινική κωδικοποίηση στο listings και αντιστοίχιση ελληνικών (pdfLaTeX)· καταχώριση ελληνικών χαρακτήρων (XeLaTeX)\\
\verb|\url| τυπώνεται «ηττπς://ωωω…»
& Κωδικοποίηση LGR
& Τα URL γράφονται πάντα με λατινικά\\
Το \texttt{;} τυπώνεται ως άνω τελεία «·»
& Σύμβαση LGR
& Ιδιότητα \texttt{keep-semicolon} του babel-greek\\
Βιβλιογραφία: «\foreignlanguage{english}{Diplwmatik'h ergas'ia}», «\foreignlanguage{english}{kai}», «\foreignlanguage{english}{Sto}» σε XeLaTeX
& Το \texttt{ntua.bst} είχε ελληνικά σε μεταγραφή LGR
& Μετατροπή σε Unicode (δουλεύει και στις τρεις μηχανές)\\
Μικρά κεφαλαία χωρίς ελληνικά
& Οι OpenType γραμματοσειρές δεν έχουν (ή έχουν σπασμένα) ελληνικά μικρά κεφαλαία
& Ελληνικά μικρά κεφαλαία χωρίς τόνους, όπως στο pdfLaTeX\\
Σφάλμα μετά από αλλαγή compiler pdfLaTeX→XeLaTeX
& Τα παλιά \texttt{.aux} περιέχουν εντολές μόνο του LGR
& Ορίζονται και στο XeLaTeX\\
\bottomrule
\end{longtable}

\section{Απόδειξη λειτουργίας}
\subsection{Το αρχικό template}
Με το ίδιο TeX Live (\texversion), ένα πέρασμα του αρχικού template έδωσε:
\begin{itemize}
	\item pdfLaTeX, όλες οι επιλογές εκτός \texttt{plain,noindex}:
		\textcolor{bad}{\textbf{\origpdferrors{} σφάλματα}} TikZ·
	\item XeLaTeX, ίδιες επιλογές: \textcolor{bad}{\textbf{\origxeerrors{} σφάλματα}} (\verb|\textgreek| κ.ά.).
\end{itemize}
Η νέα έκδοση, στις ίδιες ρυθμίσεις: \textcolor{ok}{\textbf{0 σφάλματα}}.

\subsection{Πριν και μετά}
Το ίδιο κεφάλαιο ελέγχου (\texttt{tests/edgecases.tex}) σε κάθε περίπτωση.

\evidence{\textbf{Αρχικό, pdfLaTeX, χωρίς} \enc: τα αγγλικά τυπώνονται με ελληνικά γράμματα, το \texttt{;} γίνεται «·».}{before_pdflatex_no_en.png}
\evidence{\textbf{Αρχικό, pdfLaTeX, με} \enc: το \texttt{;} γίνεται «·», το URL βγαίνει «ηττπς://…».}{before_pdflatex_text.png}
\evidence{\textbf{Νέο, XeLaTeX, χωρίς κανένα} \enc: όλα σωστά.}{after_xelatex_text.png}
\evidence{\textbf{Κώδικας, αρχικό, pdfLaTeX}: \texttt{def}~→~«δεφ», τα ελληνικά σχόλια χάνονται.}{before_pdflatex_code.png}
\evidence{\textbf{Κώδικας, νέο, pdfLaTeX}.}{after_pdflatex_code.png}
\evidence{\textbf{Κώδικας, νέο, XeLaTeX}.}{after_xelatex_code.png}
\par\medskip\noindent
\begin{minipage}{\linewidth}
\textbf{Εξώφυλλο}: pdfLaTeX (Kerkis) και XeLaTeX (Libertinus) με τα ίδια στοιχεία.\par\smallskip
\shot[0.48\linewidth]{cover_pdflatex.png}\hfill\shot[0.48\linewidth]{cover_xelatex.png}
\end{minipage}

\subsection{Αυτόματοι έλεγχοι}\label{sec:tests}
Ο \texttt{tests/run\_tests.py} αντιγράφει το template σε καθαρό φάκελο,
εφαρμόζει το σενάριο, κάνει πλήρη μεταγλώττιση με \texttt{latexmk} (μαζί με
BibTeX και ευρετήριο) και ελέγχει ότι:
\begin{itemize}
	\item δεν υπάρχει κανένα σφάλμα TeX και κανένας χαρακτήρας που λείπει από τη γραμματοσειρά·
	\item δεν υπάρχουν αόριστες αναφορές ή παραπομπές·
	\item δεν χρησιμοποιούνται bitmap γραμματοσειρές· στο XeLaTeX/LuaLaTeX δεν γίνεται αντικατάσταση γραμματοσειράς·
	\item το PDF περιέχει το αναμενόμενο κείμενο (π.χ. «Machine Learning», «def hello», το URL, «Τι είναι το Deep Learning;», «Πτυχιακή εργασία»)·
	\item το PDF \emph{δεν} περιέχει γνωστά λάθη (π.χ. «ηττπ», «δεφ», «Deep Learning·», «\foreignlanguage{english}{Diplwmatik}», «\foreignlanguage{english}{ Sto }»)·
	\item στο XeLaTeX/LuaLaTeX τα μικρά κεφαλαία του εξωφύλλου δεν έχουν γίνει πεζά χωρίς τόνους (π.χ. «Διπλωματικη Εργασια»)·
	\item με \texttt{hyperref} τα μεταδεδομένα PDF έχουν τίτλο.
\end{itemize}
Το κεφάλαιο ελέγχου έχει μικτό κείμενο, ελληνικό ερωτηματικό και άνω τελεία,
διαλυτικά, εισαγωγικά, υποσημείωση, URL, μαθηματικά, TikZ, πίνακα, εικόνα,
κώδικα Python με ελληνικά σχόλια, αλγόριθμο, ευρετήριο και ελληνικές και
αγγλικές παραπομπές.

\begin{longtable}{@{}r>{\raggedright}p{0.47\linewidth}lrrl@{}}
\toprule
\# & \textbf{Σενάριο} & \textbf{Μηχανή} & \textbf{Χρόνος} & \textbf{Σελ.} & \\
\midrule\endhead
01 & Το main.tex όπως διανέμεται (plain, noindex) & pdfLaTeX & 5.7\,s & 70 & \PASS\\
02 & Όλες οι επιλογές (TikZ, υδατογράφημα, ευρετήριο) + κεφάλαιο ελέγχου με \verb|\en{}| & pdfLaTeX & 15.2\,s & 80 & \PASS\\
03 & Όλες οι επιλογές, αγγλικά \emph{χωρίς} \verb|\en{}| & XeLaTeX & 11.5\,s & 80 & \PASS\\
04 & Όλες οι επιλογές, αγγλικά \emph{χωρίς} \verb|\en{}| & LuaLaTeX & 22.0\,s & 80 & \PASS\\
05 & Γραμματοσειρά classic & pdfLaTeX & 16.7\,s & 80 & \PASS\\
06 & Γραμματοσειρά classic, χωρίς \verb|\en{}| & XeLaTeX & 12.1\,s & 80 & \PASS\\
07 & Ασπρόμαυρη έκδοση (printer), χωρίς hyperref & XeLaTeX & 10.3\,s & 80 & \PASS\\
08 & Ελάχιστες επιλογές (μόνο modern) & pdfLaTeX & 12.6\,s & 79 & \PASS\\
09 & Δύο φοιτητές, λατινικά ονόματα, μεγάλος τίτλος, κενός υπότιτλος & XeLaTeX & 11.0\,s & 78 & \PASS\\
10 & Δύο φοιτητές, λατινικά ονόματα, μεγάλος τίτλος, κενός υπότιτλος & pdfLaTeX & 15.4\,s & 78 & \PASS\\
11 & Εικόνα εξωφύλλου σε TikZ (κώδικας με «;») & pdfLaTeX & 14.9\,s & 78 & \PASS\\
12 & Εικόνα εξωφύλλου σε TikZ (κώδικας με «;») & LuaLaTeX & 21.6\,s & 78 & \PASS\\
13 & Δική μας γραμματοσειρά (\verb|\setmainfont|, CMU Serif) & XeLaTeX & 12.0\,s & 80 & \PASS\\
14 & Αλλαγή compiler με παλιά βοηθητικά αρχεία & pdfLaTeX → XeLaTeX & 23.1\,s & 80 & \PASS\\
15 & Αλλαγή compiler με παλιά βοηθητικά αρχεία & XeLaTeX → pdfLaTeX & 22.5\,s & 80 & \PASS\\
16 & Παλιό στυλ: παντού \verb|\en{}| & XeLaTeX & 11.1\,s & 80 & \PASS\\
17 & Το main.tex όπως διανέμεται (plain, noindex) & LuaLaTeX & 9.1\,s & 70 & \PASS\\
18 & Μόνο λατινικό ονοματεπώνυμο, χωρίς εικόνα εξωφύλλου & XeLaTeX & 10.7\,s & 78 & \PASS\\
19 & Μόνο λατινικό ονοματεπώνυμο, χωρίς εικόνα εξωφύλλου & pdfLaTeX & 15.1\,s & 78 & \PASS\\
\bottomrule
\end{longtable}
\noindent Αποτέλεσμα: \textbf{19 από 19} σενάρια πέρασαν. Ο χρόνος είναι για πλήρη μεταγλώττιση από την αρχή (όλα τα περάσματα, BibTeX, ευρετήριο).


\subsection{Χρόνοι μεταγλώττισης}
Ένα πέρασμα της μηχανής για όλο το δείγμα (\pagesample{} σελίδες), με τα
βοηθητικά αρχεία ήδη έτοιμα. Το πρώτο compile χρειάζεται 3–4 περάσματα.
\begin{center}\begin{tabular}{@{}llrr@{}}
\toprule
\textbf{Μηχανή} & \textbf{Επιλογές} & \textbf{Ένα πέρασμα} & \textbf{Σελίδες}\\
\midrule
pdfLaTeX & plain,noindex & 1.5\,s & 70\\
pdfLaTeX & όλες & 3.2\,s & 78\\
XeLaTeX & plain,noindex & 1.5\,s & 70\\
XeLaTeX & όλες & 2.9\,s & 78\\
LuaLaTeX & plain,noindex & 2.5\,s & 70\\
LuaLaTeX & όλες & 4.7\,s & 78\\
\bottomrule
\end{tabular}\end{center}

Με \texttt{plain,noindex} ένα πέρασμα παίρνει περίπου 1,5\,s εδώ, οπότε
ακόμη και το πρώτο compile (3–4 περάσματα) έχει περιθώριο μέσα στο όριο των
10\,s του δωρεάν Overleaf. Με όλες τις επιλογές το πρώτο compile μπορεί να το
ξεπεράσει. Οι μετρήσεις είναι από τον υπολογιστή ελέγχου· στο Overleaf οι
χρόνοι διαφέρουν.

\section{Χρήση}
\begin{itemize}
	\item \textbf{Overleaf}: New Project → Upload Project → \texttt{overleaf-template.zip}·
		Menu → Compiler: \textbf{XeLaTeX}· TeX Live: η πιο πρόσφατη.
	\item \textbf{Τοπικά}: \verb|latexmk -xelatex main.tex| (ή \verb|-pdf|, \verb|-lualatex|).
	\item \textbf{Γραμματοσειρά} (XeLaTeX): \verb|\setmainfont{...}| μετά το \verb|\documentclass|.
	\item \textbf{Έλεγχοι}: \verb|python3 tests/run_tests.py|.
	\item \textbf{Αναφορά και σελίδα}: \verb|python3 docs/build_docs.py --run-tests|
		ξαναφτιάχνει αυτή την αναφορά, το \texttt{overleaf-template.zip} και τα
		αποτελέσματα· η σελίδα του έργου είναι το \texttt{docs/index.html}
		(GitHub Pages), με κουμπί «Άνοιγμα στο Overleaf».
\end{itemize}

\section{Γνωστοί περιορισμοί}
\begin{itemize}
	\item Με pdfLaTeX η αναζήτηση/αντιγραφή κειμένου από το PDF δίνει για
		ορισμένα γράμματα όμοιους αλλά διαφορετικούς χαρακτήρες (π.χ.
		«µ» αντί για «μ»). Το XeLaTeX δεν έχει αυτό το πρόβλημα.
	\item Με pdfLaTeX, αν ορίσετε δικό σας \texttt{literate} στο
		\verb|\lstset|, αντικαθιστάτε και την αντιστοίχιση των ελληνικών
		χαρακτήρων στον κώδικα.
	\item Το δείγμα \texttt{figures/peerArchitecture.eps} περιέχει δική του
		bitmap γραμματοσειρά· αντικαταστήστε το με τα δικά σας σχήματα.
	\item Με πλήρη γραφικά και ευρετήριο, το πρώτο compile στο δωρεάν
		Overleaf μπορεί να ξεπεράσει τα 10\,s.
\end{itemize}

\end{document}
